\documentclass[12pt]{article}
\usepackage{amsmath, amssymb}
\usepackage{geometry}
\geometry{margin=1in}

\title{CSE 400 Lecture Scribe}
\date{}

\begin{document}

\maketitle

\section*{CSE 400: Fundamentals of Probability in Computing}

\textbf{Lecture Scribe: Joint Random Variables and Joint Distributions (L15--L17)}

\textbf{Instructor:} Dr. Dhaval Patel

\textbf{Course:} CSE400 -- Fundamentals of Probability in Computing

\bigskip

\textbf{Topic:} Joint Random Variables and Joint Distributions

\textbf{Date:} March 10, 2026

\section{Motivation for Joint Random Variables}

In many real-world systems, the behavior of interest depends on multiple random variables simultaneously. Instead of analyzing a single variable, we must analyze relationships between variables.

Let

\[
X, Y
\]

be two random variables defined on the same probability space.

The goal is to understand the probability structure of the pair $(X,Y)$ rather than each variable independently.

\subsection{Example: Smart Transportation}

Consider a smart traffic monitoring system with two variables:

\[
X = \text{Vehicle Speed}
\]

\[
Y = \text{Traffic Density}
\]

Question:

Can speed be high when traffic density is also high?

This requires studying joint behavior, since the variables may be correlated.

Applications in Intelligent Transportation Systems (ITS):

\begin{itemize}
\item Adaptive traffic signal control
\item Congestion prediction
\item Autonomous vehicle routing
\end{itemize}

\subsection{Example: Wireless Network Performance}

In wireless communication systems:

\[
X = \text{Signal-to-Noise Ratio (SNR)}
\]

\[
Y = \text{Data Throughput}
\]

Question:

If SNR is high, will throughput always be high?

Answer: Not necessarily, because network congestion may affect throughput.

Applications:

\begin{itemize}
\item 5G scheduling
\item Adaptive modulation
\item Network capacity planning
\end{itemize}

\subsection{Additional Examples}

Other scenarios involving multiple random variables include:

\begin{enumerate}
\item \textbf{Weather Prediction}

\[
X = \text{Rainfall}, \quad Y = \text{Temperature}
\]

\item \textbf{Drone Delivery Systems}

\[
X = \text{Wind Speed}, \quad Y = \text{Delivery Time}
\]

\item \textbf{Rolling Two Dice}

Let

\[
X = \text{Outcome of first die}
\]

\[
Y = \text{Outcome of second die}
\]

Questions such as

\[
X + Y = 7 \quad \text{or} \quad X > Y
\]

require the joint probability mass function.

\end{enumerate}

Thus, many applications require studying multiple random variables simultaneously.

\section{Discrete Case: Joint Probability Mass Function (Joint PMF)}

\subsection{Definition}

For two discrete random variables $X$ and $Y$, the joint probability mass function (PMF) is defined as

\[
p_{X,Y}(x,y) = P(X=x,Y=y)
\]

Properties:

\begin{enumerate}
\item Non-negativity

\[
p_{X,Y}(x,y) \ge 0
\]

\item Total probability equals one

\[
\sum_x \sum_y p_{X,Y}(x,y) = 1
\]

\end{enumerate}

\subsection{Joint PMF Table Representation}

A convenient way to represent a joint PMF is using a table.

\[
\begin{array}{c|ccc}
X\backslash Y & y_1 & y_2 & \dots \\
\hline
x_1 & p(x_1,y_1) & p(x_1,y_2) & \dots \\
x_2 & p(x_2,y_1) & p(x_2,y_2) & \dots \\
\vdots & \vdots & \vdots & \dots
\end{array}
\]

Each cell contains

\[
P(X=x_i,Y=y_j)
\]

and the entire table sums to 1.

\section{Example: Rolling Two Dice}

Consider two fair dice.

Define

\[
X = \text{Outcome of the first die}
\]

\[
Y = \text{Outcome of the second die}
\]

Possible values:

\[
X,Y \in \{1,2,3,4,5,6\}
\]

Each ordered pair

\[
(x,y)
\]

occurs with probability

\[
P(X=x,Y=y) = \frac{1}{36}
\]

since there are

\[
6 \times 6 = 36
\]

equally likely outcomes.

Thus,

\[
p_{X,Y}(x,y) = \frac{1}{36}
\]

for all $x,y \in \{1,\dots,6\}$.

This joint PMF enables answering questions such as

\[
P(X+Y=7)
\]

or

\[
P(X>Y)
\]

which require considering pairs of outcomes.

\section{Continuous Case: Joint Probability Density Function}

For continuous random variables, probabilities are defined using a joint probability density function (PDF).

\subsection{Definition}

Let $X$ and $Y$ be continuous random variables.

The joint PDF is denoted

\[
f_{X,Y}(x,y)
\]

such that for any region $A$:

\[
P((X,Y)\in A) = \iint_A f_{X,Y}(x,y)\,dx\,dy
\]

\subsection{Properties}

\begin{enumerate}
\item Non-negativity

\[
f_{X,Y}(x,y) \ge 0
\]

\item Total probability equals 1

\[
\int_{-\infty}^{\infty} \int_{-\infty}^{\infty} f_{X,Y}(x,y)\,dx\,dy = 1
\]

\end{enumerate}

\subsection{Geometric Interpretation}

For continuous random variables:

\[
P(a \le X \le b, c \le Y \le d)
\]

corresponds to the volume under the joint density surface over the rectangular region

\[
[a,b] \times [c,d]
\]

Thus,

\[
P(a \le X \le b, c \le Y \le d)
=
\int_a^b \int_c^d f_{X,Y}(x,y)\,dy\,dx
\]

\section{Worked Example (Continuous Case)}

Suppose the joint PDF is defined over a specific region.

To verify validity:

\begin{enumerate}
\item Ensure

\[
f_{X,Y}(x,y) \ge 0
\]

\item Compute normalization constant $c$ using

\[
\iint f_{X,Y}(x,y)\,dx\,dy = 1
\]

\end{enumerate}

Example step:

\[
\int_0^1 \int_0^1 c(x+y)\,dx\,dy = 1
\]

Compute inner integral:

\[
\int_0^1 (x+y)\,dx
\]

Then integrate with respect to $y$, and solve for $c$.

This ensures the density integrates to 1.

\section{Joint Cumulative Distribution Function (Joint CDF)}

\subsection{Definition}

The joint cumulative distribution function of $X$ and $Y$ is

\[
F_{X,Y}(x,y) = P(X \le x, Y \le y)
\]

This function describes the probability that both random variables are simultaneously below given thresholds.

\section{Joint CDF for Continuous Variables}

For continuous random variables:

\[
F_{X,Y}(x,y)
=
\int_{-\infty}^{x} \int_{-\infty}^{y}
f_{X,Y}(u,v)\,dv\,du
\]

Thus, the joint CDF is obtained by integrating the joint PDF.

\section{Joint CDF for Discrete Variables}

For discrete random variables:

\[
F_{X,Y}(x,y)
=
\sum_{u\le x}\sum_{v\le y} p_{X,Y}(u,v)
\]

Thus the CDF is obtained by summing joint PMF values.

\section{Properties of the Joint CDF}

\subsection{Bounds}

\[
0 \le F_{X,Y}(x,y) \le 1
\]

\subsection{Monotonicity}

The function is non-decreasing in both arguments.

If

\[
x_1 \le x_2
\]

then

\[
F_{X,Y}(x_1,y) \le F_{X,Y}(x_2,y)
\]

Similarly for $y$.

\subsection{Limits}

\[
\lim_{x\to -\infty} F_{X,Y}(x,y) = 0
\]

\[
\lim_{y\to -\infty} F_{X,Y}(x,y) = 0
\]

\[
\lim_{x,y\to \infty} F_{X,Y}(x,y) = 1
\]

\subsection{Probability of a Rectangle}

For any region

\[
a < X \le b, \quad c < Y \le d
\]

the probability can be computed using the joint CDF:

\[
P(a < X \le b, c < Y \le d)
=
F(b,d) - F(a,d) - F(b,c) + F(a,c)
\]

\section{Summary}

This lecture introduced the concept of joint random variables and their distributions.

Main components studied:

\begin{enumerate}
\item Motivation for analyzing multiple random variables.
\item Joint PMF for discrete variables.
\item Joint PDF for continuous variables.
\item Geometric interpretation of joint densities.
\item Joint cumulative distribution function (CDF).
\item Properties and computation of probabilities using the joint CDF.
\end{enumerate}

These tools allow modeling interdependent stochastic phenomena, which are fundamental in:

\begin{itemize}
\item network systems
\item intelligent transportation
\item wireless communication
\item environmental modeling
\item probabilistic computing systems
\end{itemize}

\end{document}